\documentclass[12pt,a4paper]{article}

\usepackage[utf8]{inputenc}
\usepackage[T1]{fontenc}
\usepackage{lmodern}
\usepackage[margin=2.5cm]{geometry}
\usepackage{graphicx}
\usepackage{booktabs}
\usepackage{hyperref}
\usepackage{xcolor}
\usepackage{float}
\usepackage{amsmath}
\usepackage{caption}
\usepackage{subcaption}
\usepackage{parskip}
\usepackage{enumitem}
\usepackage{multirow}

\hypersetup{colorlinks=true, linkcolor=blue!60!black, citecolor=blue!60!black, urlcolor=blue!60!black}

\title{\LARGE\textbf{Do Negative Airbnb Reviews Reflect Bad Quality\\or Simply a Bad Fit?}\\[0.5em]
\large A Data-Driven Analysis of 2 Million London Reviews\\Using NLP and GPU-Accelerated Machine Learning}

\author{Abhinav Bijpuria \and Marta Trales \and Rares Olteanu}
\date{May 2026}

\begin{document}
\maketitle

% ============================================================
\begin{abstract}
Not every negative Airbnb review means something is wrong with the property.
Sometimes the listing is perfectly fine, but it just was not the right match for that particular guest.
For example, a small city-centre studio might score five stars from solo travellers but receive complaints from a family of four.

In this project we analyse \textbf{2,097,996 Airbnb reviews from London} to find out how often a negative review signals a real quality problem (like dirt or broken equipment) versus a \emph{mismatch} between the guest's expectations and what the listing actually offers (like the room being too small or the neighbourhood being too noisy for their taste).

We combine rule-based keyword detection with a \textbf{GPU-accelerated neural network} that achieves \textbf{87\% accuracy} at classifying ambiguous reviews.
Our main finding: among the 105,365 negative reviews, roughly \textbf{51\% point to genuine deficiencies} while \textbf{38\% are better explained as a mismatch}, which suggests that a large share of negative feedback could be addressed through better listing descriptions rather than property improvements.
\end{abstract}

% ============================================================
\section{Introduction}

Airbnb hosts live and die by their review scores.
A string of negative reviews can push a listing down in search results, reduce bookings, and ultimately cost the host money.
But what if many of those negative reviews are not really about the quality of the listing at all?

Consider two situations:
\begin{itemize}[nosep]
    \item \textbf{Quality deficiency:} The bathroom has mould, the WiFi does not work, and the host never replies to messages. These are real problems that need fixing.
    \item \textbf{Mismatch:} The flat is clean, well-maintained, and the host is friendly --- but a guest complains that the neighbourhood is ``too far from the centre'' or the room is ``too small.'' The listing was not a good match for that guest, but there is nothing actually wrong with it.
\end{itemize}

Telling these two cases apart matters.
If a host is getting complaints about dirt, they need to clean better.
If guests keep saying the room is small, the host might just need to add clearer photos and a note about the room size in the description.

In this report we build a system that reads every negative review, figures out which aspects the guest is complaining about, checks the listing's overall quality scores, and assigns a label: \emph{deficiency}, \emph{mismatch}, \emph{mixed}, or \emph{unclear}.
For the many reviews that are hard to classify with simple rules, we train a neural network on the GPU to make a final decision.

% ============================================================
\section{Data}

We use the publicly available \textbf{Inside Airbnb} dataset for London, scraped in September 2025.
The main file (\texttt{notes.csv}) is a pre-merged table that joins every review with its parent listing's metadata.

\begin{table}[H]
\centering
\caption{Dataset overview}
\begin{tabular}{lr}
\toprule
\textbf{Metric} & \textbf{Value} \\
\midrule
Total reviews          & 2,097,996 \\
Unique listings        & 32,021 \\
Negative reviews (VADER $\leq$ --0.05) & 105,365\,(5.0\%) \\
File size              & 6\,GB (89 columns) \\
\bottomrule
\end{tabular}
\end{table}

Each row contains the review text, reviewer ID, listing ID, room type, price, host response rate, superhost status, and six review-score dimensions (rating, accuracy, cleanliness, communication, location, value).

% ============================================================
\section{Method}

Our pipeline has four main stages.

\subsection{Stage 1 --- Sentiment Scoring with VADER}

We first need to decide which reviews are negative.
We use \textbf{VADER} (Valence Aware Dictionary and sEntiment Reasoner), a well-known tool that gives every piece of text a \emph{compound score} between --1 (very negative) and +1 (very positive).

Any review with a compound score of \textbf{--0.05 or below} is treated as negative.
This threshold is intentionally lenient so that we catch mildly negative reviews too, not just strongly angry ones.

Out of nearly 2.1 million reviews, \textbf{105,365 (5.0\%)} are flagged as negative.

\subsection{Stage 2 --- Aspect Tagging with Keyword Patterns}

Next, we identify \emph{what} the guest is complaining about.
We define eight aspect categories, each with a list of keywords, and scan every negative review for matches.

\begin{table}[H]
\centering
\caption{Aspect categories and example keywords}
\begin{tabular}{llp{7cm}}
\toprule
\textbf{Type} & \textbf{Aspect} & \textbf{Example keywords} \\
\midrule
\multirow{4}{*}{\rotatebox{90}{\small Quality}}
& Cleanliness     & dirty, filthy, mould, smell, stain, bugs \\
& Maintenance     & broken, leak, heating, WiFi, bedbug, repair \\
& Host service    & rude, unresponsive, late, cancelled, keys \\
& Accuracy/Safety & misleading, not as described, unsafe, scam \\
\midrule
\multirow{4}{*}{\rotatebox{90}{\small Mismatch}}
& Location fit    & far, commute, tube, central, zone, area \\
& Space fit       & small, tiny, cramped, narrow, stairs \\
& Value fit       & expensive, overpriced, not worth, cost \\
& Preference fit  & noisy, loud, hot, cold, dark, firm bed \\
\bottomrule
\end{tabular}
\end{table}

We use regex word-boundary matching (\verb|\b|) with non-capturing groups so that ``location'' matches but ``relocation'' does not.

\subsection{Stage 3 --- Rule-Based Classification}

For each negative review, we compute two scores:

\begin{itemize}[nosep]
    \item A \textbf{deficiency score} that goes up when quality-related keywords appear, the listing has low review scores in the relevant dimension, or similar complaints keep showing up from other guests at the same listing.
    \item A \textbf{mismatch score} that goes up when preference/fit keywords appear but the listing actually has \emph{high} quality scores (meaning other guests were happy --- this guest just was not the right audience).
\end{itemize}

The label is then assigned by comparing the two scores:

\begin{table}[H]
\centering
\caption{Labelling rules}
\begin{tabular}{lp{9cm}}
\toprule
\textbf{Label} & \textbf{Rule} \\
\midrule
Deficiency & Deficiency score $\geq$ mismatch score $+$ 1 \\
Mismatch   & Mismatch score $\geq$ deficiency score $+$ 1 \\
Mixed      & Both scores $>$ 0 but neither dominates \\
Unclear    & Negative review but no strong signal either way \\
\bottomrule
\end{tabular}
\end{table}

At this stage, a large share of reviews (\textbf{68.3\%}) land in the ``unclear'' bucket --- they are negative but do not contain enough keywords for a confident rule-based decision.
This is where the neural network comes in.

\subsection{Stage 4 --- GPU-Accelerated Neural Text Classifier}

To resolve the unclear reviews, we train a \textbf{PyTorch multi-layer perceptron (MLP)} on the GPU (NVIDIA GeForce RTX 3050).

\begin{enumerate}[nosep]
    \item \textbf{Feature extraction:} We convert each review into a numerical vector using TF-IDF (Term Frequency--Inverse Document Frequency) with up to 10,000 features and bigrams.
    \item \textbf{Training data:} The 29,282 reviews that the rule-based step labelled as either ``deficiency'' or ``mismatch'' serve as the training set (80/20 train/test split).
    \item \textbf{Architecture:} A three-layer MLP (10,000 $\rightarrow$ 256 $\rightarrow$ 128 $\rightarrow$ 2) with ReLU activations and dropout, trained for 15 epochs with Adam optimiser.
    \item \textbf{Prediction:} The model predicts a probability for each unclear review. If the model is at least \textbf{70\% confident}, we accept its label; otherwise the review stays ``unclear.''
\end{enumerate}

\begin{table}[H]
\centering
\caption{GPU model performance on the held-out test set}
\begin{tabular}{lccc}
\toprule
\textbf{Class} & \textbf{Precision} & \textbf{Recall} & \textbf{F1-Score} \\
\midrule
Deficiency & 0.908 & 0.893 & 0.901 \\
Mismatch   & 0.797 & 0.823 & 0.810 \\
\midrule
\textbf{Overall accuracy} & \multicolumn{3}{c}{\textbf{86.96\%}} \\
\bottomrule
\end{tabular}
\end{table}

The model successfully relabelled \textbf{64,602 out of 71,999} unclear reviews (89.7\%), leaving only 7.0\% still unresolved.

% ============================================================
\section{Results}

\subsection{Label Distribution}

After combining rule-based labels with the GPU model's predictions, the final distribution of the 105,365 negative reviews is:

\begin{table}[H]
\centering
\caption{Final (expanded) label distribution}
\begin{tabular}{lrr}
\toprule
\textbf{Label} & \textbf{Count} & \textbf{Share} \\
\midrule
Deficiency & 54,024 & 51.27\% \\
Mismatch   & 39,860 & 37.83\% \\
Mixed      & 4,084  & 3.88\% \\
Unclear    & 7,397  & 7.02\% \\
\bottomrule
\end{tabular}
\end{table}

The key takeaway: \textbf{nearly 38\% of all negative reviews are better explained as a mismatch} rather than a quality problem.
This is a substantial share that hosts could address through better listing descriptions, clearer photos, or setting guest expectations upfront.

\begin{figure}[H]
\centering
\includegraphics[width=0.75\textwidth]{figures/label_distribution_comparison.png}
\caption{Comparison of rule-based labels (left bars) vs.\ expanded labels after the GPU model fills in unclear cases (right bars). The model dramatically reduces the ``unclear'' category.}
\end{figure}

\subsection{Labels by Room Type}

\begin{figure}[H]
\centering
\includegraphics[width=0.75\textwidth]{figures/room_type_label_mix.png}
\caption{Stacked label proportions by room type. Hotel rooms have the highest deficiency share, while private rooms show a higher mismatch proportion.}
\end{figure}

\begin{table}[H]
\centering
\caption{Negative review counts by room type and expanded label}
\begin{tabular}{lrrrr}
\toprule
\textbf{Room type} & \textbf{Deficiency} & \textbf{Mismatch} & \textbf{Mixed} & \textbf{Unclear} \\
\midrule
Entire home/apt & 32,079 & 22,127 & 2,752 & 4,090 \\
Private room    & 21,466 & 17,451 & 1,297 & 3,252 \\
Hotel room      & 175    & 103    & 19    & 18 \\
Shared room     & 304    & 179    & 16    & 37 \\
\bottomrule
\end{tabular}
\end{table}

\subsection{Aspect Analysis}

\begin{figure}[H]
\centering
\includegraphics[width=0.85\textwidth]{figures/aspect_heatmap.png}
\caption{Heatmap showing how often each aspect keyword appears in reviews of each label type. ``Mixed'' reviews have the highest keyword density across all aspects.}
\end{figure}

The heatmap confirms our intuition:
\begin{itemize}[nosep]
    \item \textbf{Deficiency} reviews talk most about \emph{host service} (20.2\%), \emph{maintenance} (14.4\%), and \emph{cleanliness} (11.7\%).
    \item \textbf{Mismatch} reviews focus on \emph{location fit} (17.2\%), \emph{preference fit} (9.4\%), and \emph{space fit} (7.9\%).
    \item \textbf{Mixed} reviews mention both types heavily, with \emph{location fit} at 72.2\% and \emph{maintenance} at 41.9\%.
\end{itemize}

\subsection{Listing Quality by Label}

\begin{figure}[H]
\centering
\includegraphics[width=0.65\textwidth]{figures/avg_listing_rating_by_label.png}
\caption{Listings associated with ``mismatch'' reviews tend to have \emph{higher} average ratings than those with ``deficiency'' reviews, supporting the interpretation that these are good properties that simply did not suit the guest.}
\end{figure}

This is an important validation: listings labelled as ``mismatch'' have higher average ratings than those labelled as ``deficiency.''
This makes sense --- a listing with consistently high ratings that still gets the occasional negative review is more likely experiencing a fit issue than a quality issue.

\subsection{Word Analysis}

\begin{figure}[H]
\centering
\includegraphics[width=\textwidth]{figures/word_clouds.png}
\caption{Word clouds for all reviews (left), negative reviews (centre), and non-negative reviews (right).}
\end{figure}

\begin{figure}[H]
\centering
\includegraphics[width=0.85\textwidth]{figures/top_words_positive_negative.png}
\caption{Top filtered tokens in negative vs.\ non-negative reviews. Negative reviews prominently feature words like ``check,'' ``night,'' ``time,'' and ``bathroom,'' while positive reviews feature ``clean,'' ``location,'' and ``recommend.''}
\end{figure}

% ============================================================
\section{Technical Setup}

All computation was performed on a single laptop:

\begin{table}[H]
\centering
\caption{Hardware and software environment}
\begin{tabular}{ll}
\toprule
\textbf{Component} & \textbf{Details} \\
\midrule
GPU     & NVIDIA GeForce RTX 3050 Laptop (4\,GB VRAM) \\
CUDA    & 12.4 (via PyTorch) \\
cuDNN   & 9.1 \\
Python  & 3.11 (Miniconda) \\
PyTorch & 2.6.0+cu124 \\
Key libraries & scikit-learn, VADER, pandas, matplotlib \\
\bottomrule
\end{tabular}
\end{table}

The GPU was used for training and inference of the neural text classifier.
The entire pipeline --- loading 6\,GB of data, scoring 2.1 million reviews with VADER, tagging aspects, running the rule-based classifier, and training the neural network --- completes in approximately \textbf{22 minutes} on this hardware.

% ============================================================
\section{Limitations}

\begin{itemize}[nosep]
    \item \textbf{VADER is rule-based.} It works well for English but misses sarcasm and performs poorly on non-English reviews (the London dataset contains reviews in French, German, Spanish, and other languages).
    \item \textbf{Keyword lists are hand-crafted.} We may miss relevant aspects or include terms that are too broad (e.g.\ ``hot'' could mean weather or a compliment).
    \item \textbf{The neural model trains on rule-based labels.} Since the training data comes from our own rules, the model inherits any biases in those rules.
    \item \textbf{4\,GB VRAM limits model size.} With more GPU memory, we could use transformer-based models (like BERT) which would likely improve accuracy further.
\end{itemize}

% ============================================================
\section{Conclusion}

We built a complete pipeline that reads 2.1 million Airbnb reviews, identifies the 5\% that are negative, and classifies each one as either a quality deficiency or a guest--listing mismatch.

The main finding is that \textbf{nearly 4 in 10 negative reviews} are not about the property being bad --- they are about it being the wrong fit for that guest.
This has a practical implication: hosts receiving mismatch-type complaints might benefit more from improving their listing description (being honest about room size, distance from transport, noise levels) than from expensive renovations.

For Airbnb as a platform, separating mismatch from deficiency could lead to fairer review systems that distinguish between ``this place has problems'' and ``this place was not for me.''

\end{document}
